% Chapter — U.S. Standard Atmosphere 1976 (FR-8). Follows the mandatory
% FR-29 six-section template: purpose; assumptions and domain of validity;
% derivation; discretization and implementation; validation evidence;
% references. The equation labels eq:ussa76:geopotential, eq:ussa76:hydro,
% eq:ussa76:tm, eq:ussa76:gradient, eq:ussa76:isothermal, eq:ussa76:density,
% eq:ussa76:speedofsound, and eq:ussa76:interp are echoed verbatim in
% comments in cpp/src/models/atmosphere_ussa76.cpp (FR-29
% equation-label-to-code traceability); renaming them breaks that trace.
\chapter{Earth Atmosphere: U.S. Standard Atmosphere 1976}
\label{ch:ussa76}

\section{Purpose}
\label{sec:ussa76:purpose}

This chapter documents the Earth launch atmosphere (FR-8): the
U.S.~Standard Atmosphere, 1976 (USSA76)~\cite{ussa1976}. Below
\qty{86}{\kilo\metre} the model is fully analytic from the document's
defining constants and returns temperature, pressure, density, and speed of
sound --- the quantities the Phase~4 aerodynamics model needs for Mach
number and dynamic pressure and the propulsion model needs for
back-pressure thrust. From \qty{86}{\kilo\metre} to \qty{1000}{\kilo\metre}
the model returns density only, by documented interpolation of the
published USSA76 tables; this extension exists for profile continuity
(launch aerodynamics live below \qty{86}{\kilo\metre}, and orbit-decay drag
uses the Harris--Priester model of Chapter~\ref{ch:harrispriester}). The
model is implemented in \texttt{cpp/src/models/atmosphere\_ussa76.cpp} with
its interface in \texttt{cpp/include/star/models/atmosphere\_ussa76.hpp}.

\section{Assumptions and domain of validity}
\label{sec:ussa76:domain}

Assumptions:
\begin{enumerate}
  \item \textbf{Standard static atmosphere.} USSA76 is an idealized,
    steady-state, mid-latitude annual-mean profile~\cite{ussa1976}. There
    is no dependence on season, latitude, local time, or space weather;
    day-to-day density at altitude can differ from the standard by large
    factors. The model is the conventional reference for launch
    aerodynamics, not a weather model.
  \item \textbf{Molecular-scale temperature identification below
    \qty{86}{\kilo\metre}.} The model computes and returns the
    molecular-scale temperature $T_M$ of equation~\eqref{eq:ussa76:tm}.
    The document defines the kinetic temperature $T = T_M \cdot M/M_0$
    with a molecular-weight ratio that departs from unity only between 80
    and \qty{86}{\kilo\metre}, by at most $0.04\,\%$, and its own main
    tables print $T = T_M$ throughout that band (the correction is left
    to the user; \cite[Part~1, \S1.2.4, p.~9]{ussa1976}). This model
    adopts the same tabulation practice, so its temperature reproduces
    the printed Table~I values exactly and is at most $0.04\,\%$ above
    the rigorously corrected kinetic temperature between 80 and
    \qty{86}{\kilo\metre}. Pressure and density are unaffected (both are
    defined in terms of $T_M$).
  \item \textbf{Constant sea-level composition below
    \qty{86}{\kilo\metre}} ($M_0 = \qty{28.9644}{\kilogram\per\kilo\mole}$),
    per the document's mixing assumption~\cite[Part~1, \S1.2.4]{ussa1976}.
  \item \textbf{Tabulated densities above \qty{86}{\kilo\metre}.} The
    document computes pressure and density above \qty{86}{\kilo\metre}
    from species number-density integrals (diffusion equations for
    N$_2$, O, O$_2$, Ar, He, H; \cite[Part~1, \S1.3.1,
    eq.~33c]{ussa1976}). Those integrals are deliberately out of scope
    here: the model instead interpolates the published Table~I densities
    on a committed node grid (equation~\eqref{eq:ussa76:interp}). Only
    density is provided above \qty{86}{\kilo\metre}; temperature,
    pressure, and speed of sound are not (the printed speed-of-sound
    tables themselves terminate at \qty{86}{\kilo\metre}
    \cite[Part~1, \S1.3.10]{ussa1976}).
\end{enumerate}

Domain of validity: geometric altitude
$z \in [\qty{-5}{\kilo\metre}, \qty{86}{\kilo\metre})$ for the full
analytic state (the table's own lower limit is \qty{-5}{\kilo\metre});
$z \in [\qty{86}{\kilo\metre}, \qty{1000}{\kilo\metre}]$ for density only.
One boundary subtlety is adopted deliberately: the document equates the top
of the seventh geopotential layer, $H_7 = 84852$~m$'$, with
\qty{86}{\kilo\metre} geometric, but the exact conversion of
$z = \qty{86000}{\metre}$ gives $H = 84852.046$~m$'$ --- the stated
equivalence is a rounded convention. The implementation branches on
geometric altitude at exactly \qty{86}{\kilo\metre} and lets the layer-7
formulas extend approximately \qty{5}{\centi\metre} of geopotential past
$H_7$, so the analytic and tabulated regions meet without a gap. The
resulting seam discrepancy is print rounding only, measured at
$1.4\times10^{-5}$ relative and gated at $2\times10^{-4}$ by
\testid{ATM-USSA76-CONTINUITY}.

Out-of-domain behavior, matching the code exactly: altitudes below
\qty{-5}{\kilo\metre} or above \qty{1000}{\kilo\metre}, and non-finite
inputs, throw \texttt{std::domain\_error}; requesting the full
thermodynamic state at or above \qty{86}{\kilo\metre} also throws
\texttt{std::domain\_error} (the quantities do not exist in this model
there, and returning a fabricated value would violate the no-silent-default
rule, DX-2). The guard is gated by \testid{ATM-USSA76-DOMAIN}.

\section{Derivation}
\label{sec:ussa76:derivation}

\subsection{Geopotential altitude}

The 0--\qty{86}{\kilo\metre} region is formulated in geopotential altitude
$H$, defined so that the altitude-dependent gravity $g(z)$ is absorbed into
the vertical coordinate: $g_0'\,dH = g\,dZ$ with the inverse-square law
$g = g_0 (r_0/(r_0+Z))^2$. Integrating gives
\cite[Part~1, eq.~18]{ussa1976}
\begin{equation}
  H \;=\; \frac{r_0\,Z}{r_0 + Z},
  \label{eq:ussa76:geopotential}
\end{equation}
with $r_0 = \qty{6356766}{\metre}$ the document's effective Earth radius
and $g_0' = \qty{9.80665}{\metre\squared\per\second\squared\per\metre}$ the
unit geopotential (both from \cite[Part~1, Table~2]{ussa1976}; the ratio
$\Gamma = g_0/g_0' = 1$~m$'$/m is applied implicitly).

\subsection{Hydrostatic closed forms below \texorpdfstring{\qty{86}{\kilo\metre}}{86 km}}

The defining equations are hydrostatic equilibrium and the perfect-gas law,
combined by the document into
\begin{equation}
  \frac{dP}{P} \;=\; -\,\frac{g_0' M_0}{R^{*}\,T_M}\,dH,
  \label{eq:ussa76:hydro}
\end{equation}
where $R^{*} = \qty{8314.32}{\newton\metre\per\kilo\mole\per\kelvin}$ and
the molecular-scale temperature profile is linear in each of the seven
layers of \cite[Part~1, Table~4]{ussa1976}:
\begin{equation}
  T_M(H) \;=\; T_{M,b} + L_{M,b}\,(H - H_b),
  \qquad H_b \le H \le H_{b+1},
  \label{eq:ussa76:tm}
\end{equation}
with layer bases $H_b = \{0, 11, 20, 32, 47, 51, 71\}$~km$'$ (top
$H_7 = 84.852$~km$'$), gradients
$L_{M,b} = \{-6.5, 0, +1.0, +2.8, 0, -2.8, -2.0\}$~K/km$'$, and
$T_{M,0} = T_0 = \qty{288.15}{\kelvin}$,
$P_0 = \qty{101325}{\pascal}$.

For a gradient layer ($L_{M,b} \neq 0$), substituting
equation~\eqref{eq:ussa76:tm} into equation~\eqref{eq:ussa76:hydro} with
$dH = dT_M/L_{M,b}$ and integrating from the layer base,
\begin{equation}
  \int_{P_b}^{P} \frac{dP}{P}
   = -\,\frac{g_0' M_0}{R^{*} L_{M,b}}
     \int_{T_{M,b}}^{T_M} \frac{dT_M}{T_M}
  \;\;\Longrightarrow\;\;
  P \;=\; P_b \left[\frac{T_{M,b}}{T_{M,b} + L_{M,b}(H - H_b)}\right]
          ^{\textstyle\frac{g_0' M_0}{R^{*} L_{M,b}}},
  \label{eq:ussa76:gradient}
\end{equation}
which is the document's equation~33a. For an isothermal layer
($L_{M,b} = 0$), $T_M = T_{M,b}$ is constant and the same integral gives
the document's equation~33b:
\begin{equation}
  P \;=\; P_b \exp\!\left[-\,\frac{g_0' M_0\,(H - H_b)}{R^{*}\,T_{M,b}}
          \right].
  \label{eq:ussa76:isothermal}
\end{equation}
The base pressures $P_b$ follow by applying the appropriate member of the
pair \eqref{eq:ussa76:gradient}/\eqref{eq:ussa76:isothermal} at
$H = H_{b+1}$ successively from $b = 0$
\cite[Part~1, \S1.3.1]{ussa1976}. Density and speed of sound then follow
from the perfect-gas law and the adiabatic sound-speed relation
\cite[Part~1, eqs.~42 and~50]{ussa1976}:
\begin{equation}
  \rho \;=\; \frac{P\,M_0}{R^{*}\,T_M},
  \label{eq:ussa76:density}
\end{equation}
\begin{equation}
  c_s \;=\; \sqrt{\frac{\gamma\,R^{*}\,T_M}{M_0}},
  \qquad \gamma = 1.40 \text{ (exact, \cite[Table~2]{ussa1976})}.
  \label{eq:ussa76:speedofsound}
\end{equation}

\subsection{Tabulated density, 86--1000 km}
\label{sec:ussa76:upper}

Above \qty{86}{\kilo\metre} the committed node grid
$\{(z_i, \rho_i)\}$, transcribed from \cite[Part~4, Table~I, pp.~68--73]
{ussa1976} (102 nodes; provenance in
\texttt{tests/golden/atmosphere/manifest.toml}), is interpolated
log-linearly in geometric altitude --- equivalently, each segment is an
exponential with a constant scale height fitted to its endpoints:
\begin{equation}
  \rho(z) \;=\; \rho_i
  \exp\!\left[\frac{\ln \rho_{i+1} - \ln \rho_i}{z_{i+1} - z_i}
  \,(z - z_i)\right],
  \qquad z_i \le z \le z_{i+1}.
  \label{eq:ussa76:interp}
\end{equation}
At a node the fractional offset is identically zero and the model returns
$\rho_i$ exactly. Between nodes the interpolation error is bounded by the
log-density curvature, $\tfrac{1}{8}\,h^2 \,|\!\left(\ln\rho\right)''|$ for
node spacing $h$. The grid was therefore graded to the curvature measured
from the printed table itself: \qty{1}{\kilo\metre} spacing to
\qty{100}{\kilo\metre}, \qty{2}{\kilo\metre} to \qty{160}{\kilo\metre},
\qty{5}{\kilo\metre} to \qty{300}{\kilo\metre}, and 10--25~km above, giving
an estimated between-node error below $\sim 0.1\,\%$ everywhere ---
comparable to the $\pm$~half-unit-in-the-4th-figure quantum of the printed
values themselves. The exit-criterion gate
(\testid{ATM-USSA76-ROWS}) evaluates at published rows that are members of
the committed grid, where the reproduction is exact by construction and
therefore verifies the transcription and the evaluation path; the
between-node figure above is a documented estimate, not a gated claim.

\section{Discretization and implementation}
\label{sec:ussa76:implementation}

\begin{enumerate}
  \item \textbf{Module and traceability.} The model lives in
    \texttt{cpp/src/models/atmosphere\_ussa76.cpp}; the source comments
    echo the labels \texttt{eq:ussa76:geopotential},
    \texttt{eq:ussa76:hydro}, \texttt{eq:ussa76:tm},
    \texttt{eq:ussa76:gradient}, \texttt{eq:ussa76:isothermal},
    \texttt{eq:ussa76:density}, \texttt{eq:ussa76:speedofsound}, and
    \texttt{eq:ussa76:interp} verbatim.
  \item \textbf{Defining constants at the definition.} $T_0$, $P_0$,
    $g_0'$, $R^{*}$, $M_0$, $r_0$, $\gamma$, and the Table~4 layer arrays
    are USSA76-defining values, private to this model, and cited at their
    definitions in the source. They are deliberately \emph{not} promoted
    to \texttt{star/constants.hpp}: they define this document's
    atmosphere, not the project's physical-constant set.
  \item \textbf{Fixed-order layer walk.} Layer base values $(T_{M,b},
    P_b)$ are computed once, in fixed order $b = 0..6$, and cached in a
    function-local static; every subsequent evaluation walks the same
    fixed sequence (D-10: fixed evaluation order, single-threaded core).
  \item \textbf{Branch at exactly \qty{86}{\kilo\metre}.} $z <
    \qty{86000}{\metre}$ selects the analytic branch (layer-7 formulas
    extended $\sim\qty{5}{\centi\metre}$ of geopotential past $H_7$, per
    Section~\ref{sec:ussa76:domain}); $z \ge \qty{86000}{\metre}$ selects
    the node-grid branch, whose first node is exactly
    \qty{86000}{\metre}.
  \item \textbf{Compiled-in node table with a committed check copy.} The
    node grid is compiled into the model source, and
    \texttt{tests/golden/atmosphere/ussa76\_upper\_nodes.toml} carries the
    same values as an independently reviewable transcription check; test
    \testid{ATM-USSA76-UPPER-NODES} requires the two copies to agree
    bit-for-bit after \texttt{strtod}. A runtime file load was rejected:
    the core never parses text (D-2).
  \item \textbf{Determinism.} All paths are straight-line binary64
    arithmetic in fixed order under the D-10 flags. The gradient-layer
    power in equation~\eqref{eq:ussa76:gradient} and the exponentials use
    libm \texttt{pow}/\texttt{exp}/\texttt{log}, so cross-platform
    results may differ in the last ulp; all committed tolerances sit
    orders of magnitude above that spread.
\end{enumerate}

\section{Validation evidence}
\label{sec:ussa76:validation}

Table~\ref{tab:ussa76:validation} lists the tests that constitute this
chapter's validation evidence. Each test identifier is machine-checked to
exist in the test tree by \texttt{scripts/lint\_chapters.py}; golden
provenance and tolerances are recorded in
\texttt{tests/golden/atmosphere/manifest.toml}.

\begin{table}[htbp]
  \centering
  \caption{Validation evidence for the USSA76 model (doctest,
    \texttt{cpp/tests/test\_atmosphere.cpp}).}
  \label{tab:ussa76:validation}
  \begin{tabular}{lp{8.6cm}}
    \toprule
    Test ID & Acceptance basis \\
    \midrule
    \testid{ATM-USSA76-ROWS} &
      Phase~3 exit criterion~4: 23 published Table~I rows spanning all
      seven layers below \qty{86}{\kilo\metre} (temperature, pressure,
      density) plus 12 rows above \qty{86}{\kilo\metre} (density), each
      reproduced at print precision --- within half a unit in the 4th
      significant figure of the printed value. \\
    \testid{ATM-USSA76-UPPER-NODES} &
      The compiled-in 86--1000~km node grid equals the committed
      transcription check copy bit-for-bit, and evaluation at every node
      altitude returns the node density exactly
      (equation~\eqref{eq:ussa76:interp} with zero offset). \\
    \testid{ATM-USSA76-CONTINUITY} &
      Analytic/table seam at \qty{86}{\kilo\metre}: the analytic density
      approaching from below agrees with the first table node to
      $2\times10^{-4}$ relative (measured $1.4\times10^{-5}$; print
      rounding of the node is the dominant term). \\
    \testid{ATM-USSA76-DOMAIN} &
      Out-of-domain behavior of Section~\ref{sec:ussa76:domain}:
      \texttt{std::domain\_error} below \qty{-5}{\kilo\metre}, above
      \qty{1000}{\kilo\metre}, for non-finite input, and for full-state
      requests at or above \qty{86}{\kilo\metre}. \\
    \bottomrule
  \end{tabular}
\end{table}

\section{References}
\label{sec:ussa76:references}

This chapter relies on the U.S.~Standard Atmosphere, 1976
document~\cite{ussa1976} for every defining constant, equation, and table
value: Part~1, Table~2 (adopted constants), Table~4 (layer bases and
gradients), \S1.2.3--\S1.2.5 (geopotential and molecular-scale
temperature), \S1.3.1 (pressure equations 33a/33b), \S1.3.4 (density,
eq.~42), \S1.3.10 (speed of sound, eq.~50), and Part~4, Table~I
(geometric-altitude main tables, document pages 51--73). Vallado
\cite{vallado2013} surveys standard-atmosphere usage in astrodynamics
practice. Full bibliographic details appear in the bibliography.
